% !TEX root = ~/Code/TN_LOWE_RECONSTRUCTION/main.tex
\section{Introduction} \label{sec:INTRODUCTION}
There is a long history of advancements in particle and astroparticle physics resulting from measuring astrophysical neutrinos with energies in the \SI{10}{\mega\electronvolt} range~\cite{Ahmad2001,Abe2016}. The sun produces a continuous flux of roughly \SI{7e10}{\per\second\per\centi\meter\squared} neutrinos as a byproduct of fusion reactions within the stellar core~\cite{Vinyoles_2018}. Despite the confirmation of solar neutrino oscillations as a solution to the \textit{Solar Neutrino Problem}~\cite{PhysRevLett.65.2233}, several questions remain~\cite{Gann2021ndb}. The Deep Underground Neutrino Experiment, as next--generation large--scale neutrino detector, will have the ability to shed light on some of them.
% Solar neutrinos have been intensely studied, beginning over 50 years ago with the Homestake experiment, which observed only 39$\%$ of the neutrinos predicted by the standard solar model (SSM). Neutrino oscillations were soon postulated as a solution but were not observed until the SNO experiment measured the $\nu_e$ and total neutrino fluxes simultaneously, demonstrating flavour transitions.

\subsection{DUNE Far Detectors}
\ac{DUNE} will deploy four far detector modules in the Sanford Underground Research Facility (SURF) to survey a wide portfolio of physics topics.\iffalse The site has \SI{4700}{\metersofwater} of overburden which significantly reduces the rate of cosmic backgrounds throughout the detector and allows searches for low-energy astroparticle measurements.  The PIP-II accelerator upgrade at the \ac{FERMILab} will give a new neutrino beam directed at SURF to study long-baseline neutrino oscillations with \ac{DUNE}.\fi The far detector will be built in two phases. The first stage will include two liquid argon (LAr) time projection chambers (TPC) with approximately \SI{14}{\kilo\tonne} of active volume, which will be based on two distinct geometries, a horizontal drift (HD) and a vertical drift (VD) design. The HD will extend \SI{14.4}{\meter} in the drift direction, aligned along the horizontal coordinate $x$, \SI{12}{\meter} along the vertical coordinate $y$, and \SI{58}{\meter} along the remaining z coordinate. This module will be separated into four distinct drift volumes along the drift direction $x$. Alternating three anode plane assemblies (APA) and two cathode plane assemblies (CPA), which produce a drift potential of \SI{180}{\kilo\volt}, giving an electric field of \SI{500}{\volt\per\centi\meter}.\iffalse Each APA is outfitted with three charge readout planes (CRP), consisting of a wire mesh, to monitor the charge deposited within the active argon volume.\fi The VD far detector module simplifies the design by dividing the active mass into just two volumes of dimensions \numproduct{13.5 x 6.5 x 60}\unit{\metre\cubed} separated horizontally by the cathode plane. Both detectors incorporate a \textbf{photon detection system} (PDS) which makes use of \ac{XA} tiles \cite{Machado_2018} distributed over the APA (for HD) and cathode plane (for VD) respectively.\iffalse \ac{DUNE}'s active material, argon, is also unique in its sensitivity to the $\nu_e$ astrophysical flux.  The $\nu_e$~CC interaction on $^{40}$Ar has long been known to have a large scattering cross--section suitable for astroparticle measurements~\cite{raghavan1986inverse}. The PDS should allow for improvements in the low energy physics range, specifically regarding supernova and solar neutrinos ($\sim\SI{10}{\mega\electronvolt}$) and allow for a fast vertex reconstruction and improved calorimetric reconstruction of events above $\SI{5}{\mega\electronvolt}$.\fi

\begin{figure}[ht]
    \centering
    \includegraphics[width=.49\linewidth]{Images/fd10kt.png}
    \includegraphics[width=.49\linewidth]{Images/fdvd10kt.png}
  \caption[width=1\linewidth]{(left) \ac{DUNE} HD far detector module. (right) \ac{DUNE} VD far detector module.\label{fig:DUNE_FAR_DETECTORS}}
\end{figure}

%The second phase of \ac{DUNE} will include at least one LAr TPC. The TPC technology allows sub-cm tracking of events over a large volume with dynamic range suitable to precisely reconstruct neutrinos at several GeV for long-baseline neutrino oscillation measurements and astrophysicalñ neutrinos in the 5 to 50~MeV energy range. Several ideas for \ac{DUNE} phase2 have been proposed with potential to significantly lower thresholds, increasing reach of low-energy physics.  A near detector will also be constructed to mitigate systematic uncertainties related to long-baseline neutrino oscillation measurements.

\subsection{Background Model\label{sec:BACKGROUND_MODEL}}
The measurement of solar neutrinos in the low-energy regime (\SI{<20}{\mega\electronvolt}), particularly those from the \( ^8\mathrm{B} \), requires a thorough understanding of radiological and cosmogenic backgrounds. Radon (\(^{222}\text{Rn}\)) is a noble gas produced in the decay chain of \(^{238}\text{U}\) and is a well-known source of background in low-energy neutrino detectors~\cite{nndc_rn222,borexino_radon}. As a noble gas, it can diffuse into the active volume of liquid argon detectors, where its decay daughters may plate out on surfaces or remain suspended, producing problematic beta and gamma radiation. Among its progeny, several isotopes are particularly important for background modelling:

\begin{itemize}
  \item \textbf{\(^{214}\text{Bi}\)}: A short-lived beta emitter with a high Q-value (up to 3.27~MeV) and multiple gamma lines, making it the most significant source of prompt high-energy background~\cite{pdg2024}.
  \item \textbf{\(^{210}\text{Pb}\)}: A long-lived isotope (\(t_{1/2} = 22.3\)~years) that accumulates on detector surfaces, acting as a persistent low-energy background source~\cite{exo_radon}.
  \item \textbf{\(^{210}\text{Bi}\)}: A medium-energy beta emitter contributing a continuum of background up to 1.16~MeV.
  \item \textbf{\(^{210}\text{Po}\)}: An alpha emitter with a half-life of 138~days, often used as a surface contamination proxy due to its distinct alpha signal~\cite{saldanha_thesis}.
\end{itemize}

Understanding and mitigating the contributions from these isotopes is crucial in achieving the radiopurity required for low--energy neutrino detection. Apart from the previously mentioned isotopes, a comprehensive list of additional background sources has been considered for \ac{DUNE}'s simulation. These backgrounds can be classified into three main categories: external backgrounds, component-induced (material) backgrounds and intrinsic backgrounds from the liquid argon. Evaluating the full scope of \ac{DUNE}'s background model is out of the scope of this work. Here we focus on giving a brief summary of the main component and discussing the features that have the biggest impact on the low--energy neutrino reconstruction.

\begin{itemize}
  \item\textbf{External Backgrounds}: External radiation originates outside the detector cryostat. The flux is divided into two kinds of particles that manage to penetrate into the detector:

        \begin{enumerate}
          \item\textbf{Cavern Neutrons}: Produced by (\(\alpha,n\)) reactions in surrounding rock or via cosmic-ray spallation; they can enter the active volume and undergo elastic scattering or capture, producing a gamma--cascade which can mimic the topology of CC solar neutrino interactions. Some neutrons are captured on the cryostat components before entering the \ac{LAr}, contributing to the absolute environmental gamma flux.
          \item\textbf{Environmental gamma rays}: High-energy photons from the rock or infrastructure surrounding the detector, typically originating from natural radioactivity. The biggest contributor to this flux in \ac{DUNE}'s caver at SURF is the shotcrete covering the walls. Low--energy gammas (\SI{<2.7}{\mega\electronvolt}) entering the active volume have a huge flux contributing mostly to worsen the signal reconstruction by mimicking the de--excitation gammas of the CC interaction or making the topological distinction between \ac{ES} and \ac{CC} less reliable. The higher energy external gammas (\num{2.7} $<$E \SI{<14}{\mega\electronvolt}) can scatter and make pair production, which can mimic the complete signal topology, constituting an event candidate for some low energy analyses.
          % \item \textbf{Cosmic rays}: Even in underground environments, muons and cosmogenically produced secondaries can lead to spallation products and radiogenic activation of detector materials.
        \end{enumerate}

  \item \textbf{Component-Induced Backgrounds}: Backgrounds can arise from the materials used to construct the detector itself. Material screening and selection, as well as radon mitigation systems, are essential to reduce these sources. These include decay chains of long--lived isotopes introduced during fabrication or operation. Some of the most prominent are:

        \begin{enumerate}
          \item \textbf{Uranium (\(^{238}\)U)}: Produces alpha, beta, and gamma emissions, as well as daughters like \(^{214}\)Bi and \(^{210}\)Pb.
          \item \textbf{Thorium (\(^{232}\)Th)}: Similar to the uranium chain, with prominent gamma emitters like \(^{208}\)Tl (\SI{2.6}{\mega\electronvolt}).
          \item \textbf{Potassium (\(^{40}\)K)}: Often present in glass, \ac{PMT}s, and plastics; emits a 1.46~MeV gamma ray.
          \item \textbf{Radon (\(^{222}\)Rn)}: Present in atmospheric air and therefore very difficult to prevent.
        \end{enumerate}

  \item \textbf{Intrinsic Liquid Argon Backgrounds}: Certain radioactive isotopes are either present naturally in atmospheric argon or may be introduced as impurities. These form the intrinsic background component of \ac{LAr}:

        \begin{enumerate}
          \item \textbf{\(^{39}\)Ar}: A beta emitter with an endpoint energy of \SI{565}{\kilo\electronvolt} and a half--life of 269 years. It dominates the low--energy spectrum in atmospheric argon with an activity of \SI{\sim1}{\becquerel\per\kilo\gram}. Together with the external gammas, this background contributes to the smearing of any signal reconstruction.
          \item \textbf{\(^{42}\)Ar and \(^{42}\)K}: Rare isotopes that decay via beta emission. \(^{42}\)K emits a high-energy beta (\(Q \sim 3.5\)~MeV), which is particularly problematic due to its similarity to solar neutrino events.
          \item \textbf{\(^{85}\)Kr}: A beta emitter with a \(Q\)-value of \SI{687}{\kilo\electronvolt}, introduced via contamination from air or during purification. Often requires cryogenic distillation to remove.
          \item \textbf{Radon (\(^{222}\)Rn)}: In addition to emanation from components, radon can dissolve in liquid argon.
        \end{enumerate}
\end{itemize}

%The relative importance of these isotopes depends on the source of argon: underground argon (UAr) has significantly reduced \(^{39}\)Ar activity compared to atmospheric argon (AAr).
In \ac{DUNE}, \ac{LAr} purification and recirculation installations are planned to ensure constant operating conditions. A detailed background model informed by in situ measurements, simulations, and radiopurity assays is essential to interpret the observed data with the required precision for low--energy neutrino physics.

\begin{figure}[ht]
  \begin{subfigure}[t]{0.49\textwidth}
    \centering
    \includegraphics[width=\linewidth]{Images/ParticleID_KE_Distribution.png}
  \end{subfigure}
  \hfill
  \begin{subfigure}[t]{0.49\textwidth}
    \centering
    \includegraphics[width=\linewidth]{Images/ParticlePDG_KE_Distribution.png}
  \end{subfigure}
  \caption[width=1\linewidth]{Simulated particle flux from background model v.3 on top of predicted solar spectrum. (left) Divided by the background component. (right) Divided by particle type.\label{fig:BACKGROUND_MODEL}}
\end{figure}